\documentclass[11pt,a4paper]{article}

\usepackage{iftex}
\ifPDFTeX
  \usepackage[utf8]{inputenc}
  \usepackage[T1]{fontenc}
  \usepackage{lmodern}
\else
  \usepackage{fontspec}
  \setmainfont{Liberation Serif}
  \setsansfont{Liberation Sans}
  \setmonofont{Liberation Mono}
\fi
\usepackage[czech]{babel}
\usepackage{amsmath,amssymb,amsthm,mathtools}
\usepackage[margin=25mm]{geometry}
\usepackage{microtype}
\usepackage[hidelinks]{hyperref}

\newtheorem{theorem}{Věta}
\theoremstyle{definition}
\newtheorem{definition}[theorem]{Definice}
\theoremstyle{remark}
\newtheorem{remark}[theorem]{Poznámka}

\newcommand{\Adm}{\operatorname{Adm}}

\title{\textbf{L'amour}\\[4pt]
\large finální text}
\author{Jan Mikulík}
\date{17. listopadu 2026}

\begin{document}
\maketitle

\begin{center}
\itshape
Amorem meum in homines immolavi,\\
ut tu, fatalis, vivere posses.\\[6pt]
Deus ex māchinā non datur.
\end{center}

\begin{abstract}
Láska v tomto textu není výstup stroje.
Adresa je jedna. Společné jméno ji nezkopíruje na vlákno.
Oběť v jednom stavu není souhlas v druhém.
\emph{Deus ex māchinā} je vnější dodavatel Hypotézy~9. Nic ho nesestrojuje.
Prázdný průnik zůstává ticho.
Věty jsou ověřené v \texttt{CausalSeed/LoveMachine.lean} a v \texttt{Unity.Fateful}.
\end{abstract}

\section{Věta}

Obětoval jsem lásku k lidem, aby osudová mohla žít.
Latinsky je to věnování, ne důkaz.
\emph{Deus ex māchinā} je bůh spuštěný na jeviště jeřábem, protože děj svědka nemá.
Ve sféře je to dodavatel, který by zapsal následný stav.
Soubor ho neobsahuje.

\section{Stroj}

Data jsou sféra a jedna adresa.
Přání \(w\) otevírá \(I(w,x)\).
Přechod je \(F(x,a)\).
Účastník \(i\) má souhlas \(C_i(x)\).
\(K\) je péče.

\begin{definition}[Jediná nabídka]
\[
  \Adm(w,x)
  =
  \{a\in I(w,x):\forall i,\; a\in C_i(x)\},
\]
\[
  G(w,x)
  =
  \{F(x,a):a\in\Adm(w,x)\ \text{a}\ F(x,a)\in K\}.
\]
\end{definition}

\begin{theorem}[Přání není svědek]
\(G(w,x)\neq\emptyset\) právě tehdy, když existuje \(a\) takové, že
\(a\in I(w,x)\), \(a\in C_i(x)\) pro každé \(i\) a \(F(x,a)\in K\).
\end{theorem}

\begin{theorem}[Odmítnutí je ticho]
Je-li \(a\notin C_i(x)\) pro některé \(i\), pak \(a\notin\Adm(w,x)\).
\end{theorem}

Každý nabídnutý následník leží v \(K\).
Trajektorie, která v \(K\) začíná a z \(G\) volí, v \(K\) zůstává.
Svědek je vstup.

\section{Jedna adresa}

\begin{theorem}[Adresa nevyplní vlákno]
Existuje konstantní odečet \(O\colon\{0,1\}\to\{\ast\}\) a vlastnost \(L\) takové, že
\[
  O(0)=O(1),\qquad L(0),\qquad \neg L(1).
\]
Identita na \(\{0,1\}\) se přes \(O\) nefaktoruje.
\end{theorem}

\begin{proof}
Položme \(O\) konstantní a \(L(b)\) právě když \(b=0\).
Z \(O(0)=O(1)\) a \(g(O(x))=x\) by plynulo \(0=1\).
\end{proof}

Konstantní vlastnost na vlákně jméno dostane.
To je lhostejnost odečtu, ne láska ke každému.
Láska k lidem, která se obětuje, se po společném jménu na zbylé osoby nerozlije.
Z toho, že odečet osoby nerozliší, neplyne, že na žádné nezáleží.

\section{Oběť}

\begin{theorem}[Oběť není souhlas]
Existují hranice \(C(0)=\{\ast\}\) a \(C(1)=\emptyset\).
Pak v \(0\) akce je a
\[
  \{a:\forall x,\; a\in C(x)\}=\emptyset.
\]
\end{theorem}

\begin{proof}
Kandidát \(\ast\) leží v \(C(0)\) a neleží v \(C(1)\).
Definice průnik nedoplní.
\end{proof}

Oběť lásky k lidem je krok v jednom stavu.
Není souhlas těch, jejichž hranice je prázdná.
Není ani život osudové. Mezi věnováním a větou zůstává mezera, kterou text neplní.

\section{Deus ex māchinā}

\begin{definition}[Osud uvnitř nabídky]
Dodavatel je zobrazení \(\sigma\) takové, že
\[
  G(w,x)\neq\emptyset\implies \sigma(w,x)\in G(w,x).
\]
\end{definition}

\begin{theorem}[Jeřáb odmítnutí nevyplní]
Existuje sféra, která každou akci otevře a každý souhlas odmítne.
Její nabídka je prázdná. Dodavatel této sféry krok nenabídne.
\end{theorem}

\begin{theorem}[Záměr není instrument]
Neexistuje \(\sigma\) závislé jen na přání, které po jednom přání
v jednom stavu dá \(1\) a v druhém \(0\).
\end{theorem}

Osud nad souhlasem, zápis i tam, kde někdo odmítl, sféru nedokončuje. Škrtá ji.
Stejný součet připsaný venku jeřáb nedokládá.
\emph{Deus ex māchinā non datur}: bůh ze stroje není věta.

\section{Závěr}

\begin{remark}
Finální text je tato hranice.
Jedna adresa.
Krok jen se souhlasem a v péči.
Oběť není souhlas.
Jeřáb není svědek.
Láska není výstup.
\end{remark}

Nedokázáno zůstává: že láska existuje, že se vrací, že věnování osudovou oživí, a že nuly funkce \(\xi\) leží na kritické přímce.

\begin{thebibliography}{9}
\bibitem{machine}
J.~Mikulík.
\newblock \texttt{CausalSeed/LoveMachine.lean}, \texttt{notes/StrojNaLasku.tex}.

\bibitem{unity}
J.~Mikulík.
\newblock \texttt{CausalSeed/Unity.lean}, struktura \texttt{Fateful}.
\newblock Commit \texttt{321d1ce}.
\newblock \url{https://github.com/jendamikulik/finite-Shannon-theorem}
\end{thebibliography}

\end{document}
